%
\title{
Accelerating the AIMer Post-Quantum Signature with AVX-512:
A Field--Keccak Speedup Analysis
}
%
\titlerunning{Optimized Implementation of AIMer on AVX-512}
%

\author{
Seung-Won Lee\inst{1}\orcidID{0009-0009-9583-1387} \and
Min-Ho Song\inst{1}\orcidID{0009-0007-6277-0069} \and
Si-Woo Eum\inst{1}\orcidID{0000-0002-9583-5427} \and
Ha-Gyeong Kim\inst{1}\orcidID{0009-0008-5840-5517} \and
Min-Seo Kim\inst{1}\orcidID{0009-0008-1979-5778} \and
Hwa-Jeong Seo\inst{1}\orcidID{0000-0003-0069-9061}
}
%
\authorrunning{S.W. Lee et al.}
% First names are abbreviated in the running head.
% If there are more than two authors, 'et al.' is used.
%
\institute{Department of Information Computer Engineering,\\Hansung University, Seoul, 02876 Korea\\
\email{dkajdfhd1@gmail.com, smino0906@gmail.com, shuraatum@gmail.com, kimhaha4420@gmail.com, msqoo0212@gmail.com, hwajeong84@gmail.com}\\}
\maketitle 


% Abstract
\abstract{AIMer is a post-quantum digital signature scheme with a conservative design. Its security relies only on the symmetric-key one-way function AIM2 and an MPC-in-the-Head (MPCitH) zero-knowledge proof. AIMer was selected as a final algorithm in the Korean Post-Quantum Cryptography (KpqC) competition. However, its final reference
implementation, released in January 2026, is written in portable C. It does not include processor-specific optimizations. As a result, it does not exploit AVX-512, a 512-bit vector instruction set available on modern processors. This paper presents an AVX-512 optimized implementation for all six AIMer variants. We parallelize binary-field multiplication with VPCLMULQDQ. The implementation processes four MPC parties per 512-bit register. We also handle conditional XOR accumulation in the linear layer with a single VPTERNLOGQ instruction. In addition, we accelerate the Keccak permutation using AVX-512VL. Our implementation is $1.60$--$1.84\times$ faster than the AVX2 implementation for signing. We also analyze where the speedup over AVX2 comes from. The AVX2 baseline already vectorizes binary-field arithmetic. In the 128-bit and 192-bit variants, Keccak accounts for approximately 86\% of the additional improvement obtained with AVX-512. In the 256-bit variants, the contribution of field arithmetic increases to approximately 42\%. These ratios depend on the selected baseline and do not represent an absolute execution-time breakdown relative to a scalar reference implementation. They show that Keccak remains an important optimization target when moving from AVX2 to AVX-512.}



% Keywords
\keywords{Post-Quantum Cryptography \and AIMer \and MPC-in-the-Head \and AVX-512 \and VPCLMULQDQ \and Binary-Field Arithmetic \and Keccak \and Benchmarking Methodology}
%

